\documentclass[UTF8]{ctexart}
\usepackage{xeCJK}
\usepackage{geometry}
\geometry{left=2cm,right=2cm,top=2.5cm,bottom=2.5cm}
\setlength{\headheight}{12.7pt}

\usepackage{titlesec}
\titleformat{\section}
  {\normalfont\large\bfseries}
  {\thesection}
  {1em}
  {}
\titlespacing*{\section}
  {0pt}
  {3.5ex plus 1ex minus .2ex}
  {2.3ex plus .2ex}

\usepackage{amsmath}
\usepackage{amssymb}
\usepackage{amsthm}
\usepackage{enumitem}
\setlist[enumerate]{itemsep=0pt,topsep=0pt,parsep=0pt,leftmargin=0pt,itemindent=2.5em}
\setlist[itemize]{itemsep=0pt,topsep=0pt,parsep=0pt,leftmargin=0pt,itemindent=2.5em}
\usepackage{fancyhdr}
\pagestyle{fancy}
\fancyhf{}
\usepackage{graphicx}
\usepackage{hyperref}
\usepackage{indentfirst}
\usepackage{lmodern}


\begin{document}
\setlength{\abovedisplayskip}{1pt}
\setlength{\belowdisplayskip}{1pt}

\section{PW的可靠性}

\hypertarget{sec:定理1.1}{}
\noindent\textbf{定理1.1 (PW证明系统的可靠性)}

对任意的断言$P,Q$和程序$S$, 如果$\vdash_{\mathrm{PW}}\{P\}\,S\,\{Q\}$,
那么$\models\{P\}\,S\,\{Q\}$.

\noindent{\kaishu 证明.}

对PW规则归纳证明.

\begin{enumerate}
    \item $\mathcal{M}[\![\mathrm{skip}]\!](P)=P$, 所以$\models
    \{P\}\,\mathrm{skip}\,\{P\}$,

    \item 如果$\sigma\models P[u:=t]$
    且$[\![u:=t]\!](\sigma)=\tau$, 那么根据替换定理有
    \[
        \tau=\sigma[u:=\sigma(t)]\models P,
    \]
    所以$\models\{\,P[u:=t]\,\}\,u:=t\,\{P\}$.

    \item 假设$\models\{P\}\,S_1\,\{R\}$, $\models\{R\}\,S_2\,\{Q\}$, 则
    \[
        \mathcal{M}[\![S_1;S_2]\!](P)=\mathcal{M}[\![S_2]\!]
        \mathcal{M}[\![S_1]\!](P)\subset\mathcal{M}[\![S_2]\!](R)\subset Q,
    \]
    所以$\models\{P\}\,S_1;S_2\,\{Q\}$.

    \item 假设$\models\{P\land B\}\,S_1\,\{Q\}$,
    $\models\{P\land\neg\,B\}\,S_2\,\{Q\}$, 则
    \[
        \mathcal{M}[\![\,\mathbf{if}\,\,B\,\,\mathbf{then}\,\,S_1\,\,
        \mathbf{else}\,\,S_2\,\,\mathbf{fi}\,]\!](P)=
        \mathcal{M}[\![S_1]\!](P\land B)\cup
        \mathcal{M}[\![S_2]\!](P\land\neg B)\subset Q,
    \]
    所以$\models\{P\}\,\mathbf{if}\,\,B\,\,\mathbf{then}\,\,S_1\,\,
    \mathbf{else}\,\,S_2\,\,\mathbf{fi}\,\{Q\}$.

    \item 假设$\models\{P\land B\}\,S\,\{P\}$,
    对$k$归纳证明$\mathcal{M}[\![A^k]\!](P)\subset P\land\neg\,B$.

    \hspace{2.17em}
    首先$\mathcal{M}[\![A^0]\!](P)=\varnothing$; 假设命题对$k$成立, 那么
    \vspace{3pt}
    \[\begin{aligned}
        \mathcal{M}[\![A^{k+1}]\!](P)=\,&
        \mathcal{M}[\![\,\mathbf{if}\,\,B\,\,\mathbf{then}\,\,S;A^k\,\,
        \mathbf{else}\,\,\mathrm{skip}\,\,\mathbf{fi}\,]\!](P)\\
        =\,&\mathcal{M}[\![S;A^k]\!](P\land B)\cup
        \mathcal{M}[\![\mathrm{skip}]\!](P\land\neg\,B)\\
        =\,&\mathcal{M}[\![A^k]\!]\mathcal{M}[\![S]\!](P\land B)
        \cup(P\land\neg\,B)\\
        \subset\,&\mathcal{M}[\![A^k]\!](P)\cup(P\land\neg\,B)
        =P\land\neg\,B,\\[3pt]
    \end{aligned}\]
    归纳成立, 所以
    \[
        \mathcal{M}[\![\,\mathbf{while}\,\,B\,\,\mathbf{do}\,\,
        S\,\,\mathbf{od}\,]\!](P)=\bigcup_{k\in\mathbb{N}}\mathcal{M}
        [\![A^k]\!](P)\subset P\land\neg\,B,\\[-5pt]
    \]
    即$\models\{P\}\,\mathbf{while}\,\,B\,\,\mathbf{do}\,\,S
    \,\,\mathbf{od}\,\{P\land\neg\,B\}$.
    
    \item 假设$P\to P_1$, $\models\{P_1\}\,S\,\{Q_1\}$且$Q_1\to Q$, 则
    \[
        \mathcal{M}[\![S]\!](P)\subset\mathcal{M}[\![S]\!](P_1)
        \subset Q_1\subset Q,
    \]
    即$\models\{P\}\,S\,\{Q\}$. \hfill $\qedsymbol$
\end{enumerate}





\newpage

\section{TW的可靠性}

\hypertarget{sec:定理2.1}{}
\noindent\textbf{定理2.1 (TW证明系统的可靠性)}

对任意的断言$P,Q$和程序$S$, 如果$\vdash_{\mathrm{TW}}[P]\,S\,[Q]$,
那么$\models[P]\,S\,[Q]$.

\noindent{\kaishu 证明.}

对TW规则归纳证明.

\begin{enumerate}
    \item $\mathcal{M}_{\mathrm{tot}}[\![\mathrm{skip}]\!](P)=P$,
    所以$\models[P]\,\mathrm{skip}\,[P]$,

    \item 如果$\sigma\models P[u:=t]$, 则根据替换定理有
    \[
        [\![u:=t]\!](\sigma)=\sigma[u:=\sigma(t)]\models P,
    \]
    所以$\models[\,P[u:=t]\,]\,u:=t\,[P]$.

    \item 假设$\models[P]\,S_1\,[R]$, $\models[R]\,S_2\,[Q]$, 则
    \[
        \mathcal{M}_{\mathrm{tot}}[\![S_1;S_2]\!](P)=
        \mathcal{M}_{\mathrm{tot}}[\![S_2]\!]
        \mathcal{M}_{\mathrm{tot}}[\![S_1]\!](P)\subset
        \mathcal{M}_{\mathrm{tot}}[\![S_2]\!](R)\subset Q,
    \]
    所以$\models[P]\,S_1;S_2\,[Q]$.

    \item 假设$\models[P\land B]\,S_1\,[Q]$,
    $\models[P\land\neg\,B]\,S_2\,[Q]$, 则
    \[
        \mathcal{M}_{\mathrm{tot}}[\![\,\mathbf{if}\,\,B\,\,\mathbf{then}\,\,S_1\,\,
        \mathbf{else}\,\,S_2\,\,\mathbf{fi}\,]\!](P)=
        \mathcal{M}_{\mathrm{tot}}[\![S_1]\!](P\land B)\cup
        \mathcal{M}_{\mathrm{tot}}[\![S_2]\!](P\land\neg B)\subset Q,
    \]
    所以$\models[P]\,\mathbf{if}\,\,B\,\,\mathbf{then}\,\,S_1\,\,
    \mathbf{else}\,\,S_2\,\,\mathbf{fi}\,[Q]$.

    \item 当$t$是整型项, $z$是整型变量, 且$z$不出现在$t,B,P,S$中时, 假设
    \[
        \models[P\land B]\,S\,[P],\quad
        \models[P\land B\land t=z]\,S\,[t<z],\quad
        P\to t\geqslant 0,
    \]
    用反证法证明任意的$\sigma\models P$在$\mathbf{while}\,\,B\,\,
    \mathbf{do}\,\,S\,\,\mathbf{od}$下终止.

    \hspace{2.17em}
    假设存在循环下不终止的$\sigma\models P$, 在其中取使得$\sigma(t)\geqslant 0$
    最小的$\sigma$, 然后定义
    \[
        \tau=\sigma[z:=\sigma(t)],
    \]
    显然$\sigma\models B$. 又因为$\sigma,\tau$在$t,B,P,S$的变元上一致,
    所以根据\href{03 操作语义#nameddest=sec:定理2.2}{重合引理},
    \[
        \tau\,\text{在循环下不终止}\,,\quad \tau\models P,\quad\tau\models B,\quad
        \tau(t)=\sigma(t)=\tau(z).
    \]
    此时
    \[
        \langle\,\mathbf{while}\,\,B\,\,\mathbf{do}\,\,S\,\,
        \mathbf{od},\tau\rangle\to\langle S;\,\mathbf{while}\,\,B
        \,\,\mathbf{do}\,\,S\,\,\mathbf{od},\tau\rangle\to^*
        \langle\,\mathbf{while}\,\,B\,\,\mathbf{do}\,\,S\,\,
        \mathbf{od},[\![S]\!](\tau)\rangle\to\cdots,\\[-1pt]
    \]
    由于$\tau\models P\land B\land t=z$,
    依假设有$[\![S]\!](\tau)\models P\land t<z$.
    所以根据\href{03 操作语义#nameddest=sec:定理2.1}{框架引理},
    \vspace{1pt}
    \[
        [\![S]\!](\tau)(t)<[\![S]\!](\tau)(z)=\tau(z)=\sigma(t),\\[1pt]
    \]
    即$[\![S]\!](\tau)$也在循环下不终止且满足$P$,
    但是$[\![S]\!](\tau)(t)<\sigma(t)$, 矛盾.

    \vspace{1pt}\hspace{2.17em}
    所以任意的$\sigma\models P$在$\mathbf{while}\,\,B\,\,
    \mathbf{do}\,\,S\,\,\mathbf{od}$下终止, 同时(在PW一节)已证
    \[
        \models\{P\land B\}\,S\,\{P\}\quad\implies\quad
        \models\{P\}\,\mathbf{while}\,\,B\,\,
        \mathbf{do}\,\,S\,\,\mathbf{od}\,\{P\land\neg\,B\},
    \]
    所以$[\![\,\mathbf{while}\,\,B\,\,\mathbf{do}\,\,S\,\,\mathbf{od}\,]\!]
    (\sigma)\models P\land\neg\,B$, 即$\models[P]\,\mathbf{while}\,\,B\,\,
    \mathbf{do}\,\,S\,\,\mathbf{od}\,[P\land\neg\,B]$.

    \item 假设$P\to P_1$, $\models[P_1]\,S\,[Q_1]$且$Q_1\to Q$, 则
    \[
        \mathcal{M}_{\mathrm{tot}}[\![S]\!](P)\subset
        \mathcal{M}_{\mathrm{tot}}[\![S]\!](P_1)\subset Q_1\subset Q,
    \]
    即$\models[P]\,S\,[Q]$. \hfill $\qedsymbol$
\end{enumerate}

\end{document}
